\documentclass[11pt,class=report,crop=false]{standalone}
\usepackage[screen]{../logic}


\begin{document}


%====================================================================
\chapitre[Preuve du théorème de complétude]{Preuve du théorème\\ de complétude}
%====================================================================

% \insertvideo{}{partie 1.1. Titre}


\objectifs{Nous prouvons le théorème de complétude.}


%%%%%%%%%%%%%%%%%%%%%%%%%%%%%%%%%%%%%%%%%%%%%%%%%%%%%%%%%%%%%%%%%%%%%
\section{Le théorème de complétude}

%--------------------------------------------------------------------
\subsection{Énoncé}

On rappelle l'énoncé qu'on a déjà bien étudié dans le chapitre \og{}Théorème de complétude\fg{} mais que l'on n'a pas encore prouvé.
\begin{theoreme}[Théorème de complétude]
	Si $\Gamma \ldoubleturn \mathcal{Q}$ alors $\Gamma \lturn \mathcal{Q}$.
\end{theoreme}
\index{theoreme completude@théorème de complétude}

Autrement dit, tout ce qui est vrai est prouvable.


%--------------------------------------------------------------------
\subsection{À prouver : \og{}Si $\mathcal{T}$ est cohérente, alors $\mathcal{T}$ admet un modèle.\fg{}}

Grâce au théorème de clôture du chapitre précédent, on ramène la preuve du théorème de complétude dans le cadre général au cadre des formules fermées.
Soit $\mathcal{T}$ une théorie, $\mathcal{Q}$ une formule fermée.

\begin{theoreme}[Théorème A]
	Si $\mathcal{T} \ldoubleturn \mathcal{Q}$, alors $\mathcal{T} \lturn \mathcal{Q}$.
\end{theoreme}

On va donc étudier le cas des théories. Ce théorème A sera obtenu à partir du théorème B suivant qui lui est équivalent :

\begin{theoreme}[Théorème B]
	Si $\mathcal{T}$ est cohérente alors $\mathcal{T}$ admet un modèle.
\end{theoreme}

Ainsi on aura l'enchaînement :

\mycenterline{Théorème B $\implies$ Théorème A $\implies$ Théorème de complétude}

C'est donc ce théorème B que l'on va démontrer dans ce chapitre, après avoir justifié comment en déduire le théorème A.

%--------------------------------------------------------------------
\subsection{Rappel du vocabulaire}

\begin{itemize}
	\item \defi{Théorie}\index{theorie@théorie} : un ensemble $\mathcal{T}$ de formules fermées.
	
	\item $\mathcal{T}$ est \defi{non cohérente} si $\mathcal{T} \lturn \lantitauto$ ; sinon $\mathcal{T}$ est \defi{cohérente}\index{theorie@théorie!coherente@cohérente}.
	
	\item $\mathcal{T}$ admet un \defi{modèle}\index{modele@modèle} s'il existe une structure $\mathcal{S}$ telle que $\ldoubleturn_{\mathcal{S}} \mathcal{T}$, c'est-à-dire une structure dans laquelle toutes les formules de $\mathcal{T}$ deviennent vraies.
\end{itemize}

On a vu au chapitre précédent la proposition suivante :
\begin{proposition}
	\label{prop:noncoherente}
	$\mathcal{T} \lturn \mathcal{Q} \quad \text{ssi} \quad \mathcal{T} \cup \{ \lnot \mathcal{Q} \} \text{ non cohérente}$.
\end{proposition}

Voici comment prouver le théorème A à partir du théorème B (on a vu la réciproque dans le chapitre précédent).



\begin{proof}
	On va donc prouver le théorème A sous sa forme contraposée :
	si $\mathcal{T} \not\lturn \mathcal{Q}$, alors $\mathcal{T} \not\ldoubleturn \mathcal{Q}$.
	
	On utilisera avec modération les négations d'opérateurs : 
	$\mathcal{T} \not\lturn \mathcal{Q}$ signifie $\text{non}(\mathcal{T} \lturn \mathcal{Q})$
	et $\mathcal{T} \not\ldoubleturn \mathcal{Q}$ signifie $\text{non}(\mathcal{T} \ldoubleturn \mathcal{Q})$.	
	
	
	Hypothèse : $\text{non}(\mathcal{T} \lturn \mathcal{Q})$. 
	But : $\text{non}(\mathcal{T} \ldoubleturn \mathcal{Q})$.
	
\begin{align*}
	& \text{non}(\mathcal{T} \lturn \mathcal{Q}) \\
	\text{donc} \quad & \mathcal{T} \cup \{ \lnot \mathcal{Q} \} \text{ est cohérente} \qquad \text{par la proposition \ref{prop:noncoherente} ci-dessus} \\
	\text{donc} \quad & \mathcal{T} \cup \{ \lnot \mathcal{Q} \} \text{ admet un modèle} \qquad \text{par le théorème B} \\
	\text{donc} \quad & \ldoubleturn_{\mathcal{S}} \mathcal{T} \cup \{ \lnot \mathcal{Q} \} \qquad \text{pour une certaine structure $\mathcal{S}$} \\
	\text{donc} \quad & \ldoubleturn_{\mathcal{S}} \mathcal{T} \ \text{ et } \ \ldoubleturn_{\mathcal{S}} \lnot\mathcal{Q} \qquad \text{car dans $\mathcal{S}$, $\mathcal{T}$ est vraie et $\lnot \mathcal{Q}$ aussi} \\
	\text{donc} \quad & \text{non}(\mathcal{T} \ldoubleturn \mathcal{Q}) \qquad \text{car pour $\mathcal{S}$, $\mathcal{T}$ est vraie et $\mathcal{Q}$ fausse}
\end{align*}
	
\end{proof}


%%%%%%%%%%%%%%%%%%%%%%%%%%%%%%%%%%%%%%%%%%%%%%%%%%%%%%%%%%%%%%%%%%%%%
\section{Théorie complète}

%--------------------------------------------------------------------
\subsection{Définition}

\begin{definition}
	Une théorie $\mathcal{T}$ est dite \defi{complète}\index{theorie@théorie!complete@complète} si :
	\begin{itemize}
		\item $\mathcal{T}$ est cohérente,
		\item et pour toute formule fermée $\mathcal{Q}$, on a : $\mathcal{T} \lturn \mathcal{Q}$ ou $\mathcal{T} \lturn \lnot \mathcal{Q}$.
	\end{itemize}
\end{definition}

Ainsi dans une théorie complète, quelle que soit la formule fermée $\mathcal{Q}$, on peut la prouver ou bien prouver son contraire, mais pas les deux en même temps. La dénomination \og{}complète\fg{} est donc bien choisie. 

%--------------------------------------------------------------------
\subsection{Satisfaction}

Voici une conséquence importante quand on se place dans le cadre d'une théorie complète : les formules sont vraies ou fausses indépendamment de la structure.

\begin{proposition}
	Soient $\mathcal{T}$ une théorie complète, $\mathcal{Q}$ une formule fermée. 
	Soient $\mathcal{S}_1$ et $\mathcal{S}_2$ deux modèles de $\mathcal{T}$. 
	Alors :
	\mycenterline{$\ldoubleturn_{\mathcal{S}_1} \mathcal{Q}$ \quad ssi \quad  $\ldoubleturn_{\mathcal{S}_2} \mathcal{Q}$.}
\end{proposition}

\begin{proof}
	Comme $\mathcal{T}$ est complète, on considère les deux cas : soit $\mathcal{T} \lturn \mathcal{Q}$, soit $\mathcal{T} \lturn \lnot \mathcal{Q}$.
	\begin{itemize}
		\item Si $\mathcal{T} \lturn \mathcal{Q}$, alors $\mathcal{T} \ldoubleturn \mathcal{Q}$ par le théorème de validité, donc $\ldoubleturn_{\mathcal{S}_1} \mathcal{Q}$ et $\ldoubleturn_{\mathcal{S}_2} \mathcal{Q}$ car $\mathcal{S}_1$ et $\mathcal{S}_2$ sont des modèles de $\mathcal{T}$, c'est-à-dire que $\mathcal{Q}$ est vraie pour $\mathcal{S}_1$ et $\mathcal{S}_2$.
		
		\item Si $\mathcal{T} \lturn \lnot \mathcal{Q}$, le même raisonnement donne $\ldoubleturn_{\mathcal{S}_1} \lnot \mathcal{Q}$ et $\ldoubleturn_{\mathcal{S}_2} \lnot \mathcal{Q}$, c'est-à-dire que $\mathcal{Q}$ est fausse dans $\mathcal{S}_1$ et $\mathcal{S}_2$.
	\end{itemize}
\end{proof}


%--------------------------------------------------------------------
\subsection{Un lemme}

Dans une théorie complète, il n'y a pas de piège. On va prouver les énoncés suivants qui pourraient sembler évidents mais nécessitent cette notion.

\begin{lemme}
	\label{lem:compbis}
	Soit $\mathcal{T}$ une théorie complète. 
	\begin{enumerate}
		\item $\mathcal{T} \lturn \lnot \mathcal{P} \quad \text{ssi} \quad \text{non}(\mathcal{T} \lturn \mathcal{P})$
		\item $\mathcal{T} \lturn \mathcal{P} \land \mathcal{Q} \quad \text{ssi} \quad (\mathcal{T} \lturn \mathcal{P} \ \text{ et } \ \mathcal{T} \lturn \mathcal{Q})$
		\item $\mathcal{T} \lturn \mathcal{P} \lor \mathcal{Q} \quad \text{ssi} \quad (\mathcal{T} \lturn \mathcal{P} \ \text{ ou } \ \mathcal{T} \lturn \mathcal{Q})$
	\end{enumerate}
\end{lemme}

\begin{proof}
	~
	\begin{enumerate}
		\item 
		\begin{itemize}
			\item Supposons $\mathcal{T} \lturn \lnot \mathcal{P}$.
			Par l'absurde, si $\mathcal{T} \lturn \mathcal{P}$, alors $\mathcal{T} \lturn (\mathcal{P} \land \lnot \mathcal{P})$, donc $\mathcal{T} \lturn \lantitauto$. Contradiction avec $\mathcal{T}$ cohérente (car $\mathcal{T}$ complète).
			Donc $\text{non}(\mathcal{T} \lturn \mathcal{P})$.
			
			\item Supposons $\text{non}(\mathcal{T} \lturn \mathcal{P})$.
			Comme $\mathcal{T}$ est complète ($\mathcal{T} \lturn \mathcal{P}$ ou $\mathcal{T} \lturn \lnot \mathcal{P}$), on a donc $\mathcal{T} \lturn \lnot \mathcal{P}$.
		\end{itemize}
		
		\item 
		\begin{itemize}
			\item Si $\mathcal{T} \lturn \mathcal{P} \land \mathcal{Q}$, alors $\mathcal{T} \lturn \mathcal{P}$ et $\mathcal{T} \lturn \mathcal{Q}$ car $\mathcal{P} \land \mathcal{Q} \lturn \mathcal{P}$ et $\mathcal{P} \land \mathcal{Q} \lturn \mathcal{Q}$.
			\item Si $\mathcal{T} \lturn \mathcal{P}$ et $\mathcal{T} \lturn \mathcal{Q}$, alors $\mathcal{T} \lturn \mathcal{P} \land \mathcal{Q}$, car $\mathcal{P}, \mathcal{Q} \lturn \mathcal{P} \land \mathcal{Q}$.
		\end{itemize}
		Remarque : ce point n'utilise pas l'hypothèse que $\mathcal{T}$ est complète.
		
		\item 
		\begin{itemize}
			\item Si ($\mathcal{T} \lturn \mathcal{P}$ ou $\mathcal{T} \lturn \mathcal{Q}$), alors $\mathcal{T} \lturn \mathcal{P} \lor \mathcal{Q}$ car $\mathcal{P} \lturn \mathcal{P} \lor \mathcal{Q}$ et $\mathcal{Q} \lturn \mathcal{P} \lor \mathcal{Q}$. (Pas besoin de $\mathcal{T}$ complète).
			
			\item Supposons $\mathcal{T} \lturn \mathcal{P} \lor \mathcal{Q}$. Supposons par l'absurde que $\mathcal{T} \not\lturn \mathcal{P}$ et $\mathcal{T} \not\lturn \mathcal{Q}$.
			Comme $\mathcal{T} \not\lturn \mathcal{P}$, on a $\mathcal{T} \lturn \lnot \mathcal{P}$ car $\mathcal{T}$ est complète (par le point 1), et de même comme $\mathcal{T} \not\lturn \mathcal{Q}$, alors $\mathcal{T} \lturn \lnot \mathcal{Q}$.
			Ainsi par le point 2, $\mathcal{T} \lturn \lnot \mathcal{P} \land \lnot \mathcal{Q}$, donc $\mathcal{T} \lturn \lnot(\mathcal{P} \lor \mathcal{Q})$.
			Mais alors, on a à la fois $\mathcal{T} \lturn \mathcal{P} \lor \mathcal{Q}$ (par l'hypothèse) et $\mathcal{T} \lturn \lnot(\mathcal{P} \lor \mathcal{Q})$, donc $\mathcal{T}$ est non cohérente. C'est une contradiction car $\mathcal{T}$ complète, donc cohérente.
		\end{itemize}
	\end{enumerate}
\end{proof}


%%%%%%%%%%%%%%%%%%%%%%%%%%%%%%%%%%%%%%%%%%%%%%%%%%%%%%%%%%%%%%%%%%%%%
\section{Construction d'un modèle à partir d'une théorie complète}
\label{sec:modele}

On commence la preuve du théorème de complétude dans sa forme du théorème B : \og{}si $\mathcal{T}$ est cohérente, alors $\mathcal{T}$ admet un modèle\fg{}. On débute cette démonstration un peu par la fin.
On va partir d'une théorie complète, avec des hypothèses supplémentaires, et on va construire un modèle. On expliquera plus tard comment obtenir ces hypothèses supplémentaires.
Il existe plusieurs preuves du théorème de complétude, nous adoptons ici celle basée sur les témoins de Henkin.

%--------------------------------------------------------------------
\subsection{Hypothèses renforcées}

Soit $\mathcal{L}$ un langage et $\mathcal{T}$ une théorie. On va supposer les propriétés suivantes pour $\mathcal{T}$ :
\begin{itemize}
	\item[(i)] $\mathcal{T}$ est cohérente.
	\item[(ii)] $\mathcal{T}$ est complète.
	\item[(iii)] $\mathcal{T}$ admet des témoins de Henkin.
\end{itemize}

Remarque : (ii) $\implies$ (i) mais on insiste sur le fait qu'à la fin on aimerait avoir la seule hypothèse (i).

\begin{definition}
	$\mathcal{T}$ admet des \defi{témoins de Henkin}\index{temoin de Henkin@témoin de Henkin} si pour toute formule $\mathcal{P}$ de $\mathcal{L}$ ayant une unique variable libre $x$, il existe une constante $c_{\mathcal{P}}$ de $\mathcal{L}$ telle que la formule suivante :
	\[ \exists x \ \mathcal{P}(x) \quad \limplies \quad \mathcal{P}(c_{\mathcal{P}}) \]
	soit une formule de $\mathcal{T}$.
\end{definition}

Remarques.
\begin{itemize}
	% \item On trouve parfois une définition un peu plus faible où on exige seulement :
	% \[
	% \mathcal{T} \quad \lturn \quad  
	% \big[ (\exists x\ \mathcal{P}(x)) \ \limplies \  \mathcal{P}(c_{\mathcal{P}})\big].
	% \]
	
	\item On pourrait aussi écrire : il existe $c \in \mathcal{L}$ tel que :
	 \[
	 \text{\og}(\exists x\ \mathcal{P}(x)) \ \limplies \  \mathcal{P}(c)\text{\fg} \ \in \ \mathcal{T},
	 \]
	 si on ne souhaite pas insister sur la dépendance de $c$ vis-à-vis de $\mathcal{P}$.
	
	\item Avec des témoins de Henkin, une formule \og{}$\exists x \ \mathcal{P}(x)$\fg{} n'est plus une formule abstraite car on peut en obtenir une réalisation concrète \og{}$\mathcal{P}(c_{\mathcal{P}})$ \fg{}, d'où le nom de \og{}témoin\fg{}.
\end{itemize}


\bigskip

Voici encore deux hypothèses supplémentaires, cette fois concernant le langage $\mathcal{L}$ :
\begin{itemize}
	\item[(iv)] $\mathcal{L}$ est un langage dénombrable (l'ensemble des symboles de $\mathcal{L}$ est dénombrable). Cela entraîne que l'ensemble des formules de $\mathcal{L}$ est aussi dénombrable (voir le chapitre \og{}Théorème de compacité\fg{}).
	\item[(v)] $\mathcal{L}$ ne contient pas le prédicat d'égalité \og{}$=$ \fg{}. Sans cette hypothèse, on augmenterait la technicité de la preuve car on devrait travailler avec des classes d'équivalence de la relation définie par $t_1 \sim t_2$ lorsque $t_1 = t_2$ est vérifié.
\end{itemize}

Malgré ces hypothèses simplificatrices, la preuve reste ardue.

%--------------------------------------------------------------------
\subsection{Construction du modèle}
\label{ssec:modele}

Pour une théorie $\mathcal{T}$ qui vérifie les hypothèses (i), (ii), (iii)
dans un langage $\mathcal{L}$ qui vérifie les hypothèses (iv), (v), nous allons construire un modèle $\mathcal{S}$ pour $\mathcal{T}$, c'est-à-dire une structure $\mathcal{S}$ telle que $\ldoubleturn_{\mathcal{S}} \mathcal{T}$.

Construire la structure $\mathcal{S}$, c'est donner :
\begin{itemize}
	\item un domaine $E$,
	\item une interprétation $c^{\mathcal{S}}$ de chaque constante $c$ de $\mathcal{L}$,
	\item une interprétation $f^{\mathcal{S}}$ de chaque fonction $f$ de $\mathcal{L}$,
	\item une interprétation $P^{\mathcal{S}}$ de chaque prédicat $P$ de $\mathcal{L}$.
\end{itemize}

Ensuite, il faudra démontrer que toute formule $\mathcal{P}$ de $\mathcal{T}$ est vraie lorsqu'elle est interprétée dans $\mathcal{S}$.

\bigskip

C'est parti ! La construction est plus ou moins automatique. On définit les objets afin d'avoir les propriétés voulues.
Voici la structure $\mathcal{S}$ :
\begin{itemize}
	\item Le domaine $E$ est l'ensemble des termes fermés de $\mathcal{L}$.
	On rappelle que les termes sont formés à partir des constantes et des variables, auxquelles on peut appliquer des fonctions de manière récursive. Ici, les termes fermés sont donc uniquement les constantes et les itérations de fonctions appliquées à des constantes.
	Exemples de termes fermés : $c_1$, $c_2$, $f(c_1)$, $g(c_1, c_2)$, $f(f(c_1))$, $g(c_1,f(c_2))$, etc.
	
	Remarque : si on avait l'égalité comme élément du langage, il faudrait prendre pour $E$ l'ensemble des classes d'équivalence des termes fermés modulo la relation d'équivalence $t_1 \sim t_2$ ssi $\mathcal{T} \lturn (t_1 = t_2)$.
	
	\item Si $c$ est un symbole de constante de $\mathcal{L}$, alors $c$ est un terme fermé donc un élément de $E$. On définit simplement $c^{\mathcal{S}} = c$.
	
	\item Soit $f$ un symbole de fonction $n$-aire de $\mathcal{L}$. Définir $f^{\mathcal{S}}$, c'est définir une \og{}vraie\fg{} fonction $f^{\mathcal{S}} : E^n \to E$.
	Soient $t_1, \ldots, t_n$ des éléments de $E$. Ce sont des termes fermés de $\mathcal{L}$, donc $f(t_1, \ldots, t_n)$ est aussi un terme fermé, donc un élément de $E$. On pose $f^{\mathcal{S}}(t_1, \ldots, t_n) = f(t_1, \ldots, t_n)$.
	
	\item Soit $P$ un prédicat $n$-aire de $\mathcal{L}$. Définir l'interprétation $P^{\mathcal{S}}$, c'est définir une fonction $P^{\mathcal{S}} : E^n \to \{V, F\}$. Voici comment attribuer une valeur vrai/faux à chaque élément $(t_1, \ldots, t_n) \in E^n$ : 
	on pose $P^{\mathcal{S}}(t_1, \ldots, t_n) = V$ si $\mathcal{T} \lturn P(t_1, \ldots, t_n)$ et on pose $P^{\mathcal{S}}(t_1, \ldots, t_n) = F$ sinon.
	
	On a ainsi par construction : $\ldoubleturn_{\mathcal{S}} P(t_1, \ldots, t_n)$ ssi $\mathcal{T} \lturn P(t_1, \ldots, t_n)$. C'est une assertion que l'on généralisera plus tard à une formule quelconque.
\end{itemize}



%--------------------------------------------------------------------
\subsection{Satisfaction}

Il reste à montrer que notre structure $\mathcal{S}$ est bien un modèle pour $\mathcal{T}$, c'est-à-dire $\ldoubleturn_{\mathcal{S}} \mathcal{T}$.
La preuve repose sur le résultat suivant :

\begin{lemme}
	\label{lem:satisfaction}
	Soit $\mathcal{Q}$ une formule fermée.
	\[ \mathcal{T} \lturn \mathcal{Q} \quad \text{ssi} \quad \ldoubleturn_{\mathcal{S}} \mathcal{Q} \]
\end{lemme}

Admettons un instant ce lemme et justifions que $\mathcal{S}$ est un modèle de $\mathcal{T}$. En effet soit $\mathcal{P}$ une formule de $\mathcal{T}$, alors bien sûr $\mathcal{T} \lturn \mathcal{P}$ donc par le lemme $\ldoubleturn_{\mathcal{S}} \mathcal{P}$. Ceci étant vrai pour toute formule de $\mathcal{T}$, on a bien $\ldoubleturn_{\mathcal{S}} \mathcal{T}$.

\begin{proof}[Démonstration du lemme]
	Pour simplifier, et sans perte de généralité, on suppose que la formule $\mathcal{Q}$ ne contient que les connecteurs $\lnot$ et $\land$ ainsi que le quantificateur $\exists$. 
	Pour les connecteurs, on renvoie au chapitre \og{}Quels connecteurs sont utiles ?\fg{} afin de se ramener à ces deux connecteurs. Si on a une formule \og{}$\forall x\ \mathcal{P}(x)$\fg{}, alors elle équivaut à \og{}$\lnot(\exists x\ \lnot\mathcal{P}(x))$ \fg{} ce qui permet de transformer tous les quantificateurs $\forall$ en des quantifications $\exists$ (avec des négations en plus).
	
	Démontrons le résultat par induction sur la construction de la formule $\mathcal{Q}$.
	
	\begin{itemize}
		\item Cas $\mathcal{Q}$ est une formule atomique. 
		
		Alors $\mathcal{Q} = P(t_1, \ldots, t_n)$ où $P$ est un prédicat et $t_1,\ldots,t_n$ des termes fermés. Par construction de $\mathcal{S}$, 
		\[
		\mathcal{T} \lturn P(t_1, \ldots, t_n) \quad \text{ ssi } \quad \ldoubleturn_{\mathcal{S}} P(t_1, \ldots, t_n),
		\]
		donc le lemme est vrai dans ce cas.
		
		\item Cas $\mathcal{Q} = \lnot \mathcal{P}$. 
		
		Par récurrence on suppose le lemme vrai pour $\mathcal{P}$ et on veut le prouver pour $\mathcal{Q} = \lnot \mathcal{P}$.
		\begin{align*}
			& \mathcal{T} \lturn \mathcal{Q} \\
			\text{ ssi } \quad & \mathcal{T} \lturn \lnot \mathcal{P} \\
			\text{ ssi } \quad & \text{non}(\mathcal{T} \lturn \mathcal{P}) \qquad \text{par le Lemme \ref{lem:compbis}, car $\mathcal{T}$ est complète}\\
			\text{ ssi } \quad & \text{non}(\ldoubleturn_{\mathcal{S}} \mathcal{P}) \qquad \text{par l'hypothèse de récurrence}\\
			\text{ssi } \quad & \ldoubleturn_{\mathcal{S}} \lnot \mathcal{P}\\
			\text{ssi } \quad & \ldoubleturn_{\mathcal{S}} \mathcal{Q}			
		\end{align*}
		
		Le point clé est le passage de la deuxième à la troisième ligne qui utilise l'hypothèse (ii) $\mathcal{T}$ complète via le Lemme \ref{lem:compbis}.
		
		\item Cas $\mathcal{Q} = \mathcal{P}_1 \land \mathcal{P}_2$.
		
		Par récurrence on suppose le lemme valide pour $\mathcal{P}_1$ et $\mathcal{P}_2$.
		\begin{align*}
			& \mathcal{T} \lturn \mathcal{Q} \\
			\text{ssi } \quad & \mathcal{T} \lturn \mathcal{P}_1 \land \mathcal{P}_2 \\
			\text{ssi } \quad & (\mathcal{T} \lturn \mathcal{P}_1 \ \text{ et } \  \mathcal{T} \lturn \mathcal{P}_2) \qquad \text{par le Lemme \ref{lem:compbis}}\\
			\text{ssi } \quad & (\ldoubleturn_{\mathcal{S}} \mathcal{P}_1 \ \text{ et } \ \ldoubleturn_{\mathcal{S}} \mathcal{P}_2) \qquad \text{par l'hypothèse de récurrence}\\
			\text{ssi } \quad & \ldoubleturn_{\mathcal{S}} \mathcal{P}_1 \land \mathcal{P}_2 \\
			\text{ssi } \quad & \ldoubleturn_{\mathcal{S}} \mathcal{Q}
		\end{align*}
		

		\item Cas $\mathcal{Q} = \exists x\ \mathcal{P}(x)$.
		
		On suppose le lemme vrai pour $\mathcal{P}$.
		\begin{itemize}
			\item Si $\mathcal{T} \lturn \mathcal{Q}$, c'est-à-dire $\mathcal{T} \lturn \exists x\ \mathcal{P}(x)$. Par l'hypothèse (iii) sur les témoins de Henkin, on sait qu'il existe $c_{\mathcal{P}} \in \mathcal{L}$ tel que \og{}$(\exists x\ \mathcal{P}(x)) \limplies \mathcal{P}(c_{\mathcal{P}})$\fg{} soit une formule de $\mathcal{T}$.
			Comme $\mathcal{T} \lturn \exists x\ \mathcal{P}(x)$ et $\mathcal{T} \lturn \big[ (\exists x\ \mathcal{P}(x)) \limplies \mathcal{P}(c_{\mathcal{P}}) \big]$, alors $\mathcal{T} \lturn \mathcal{P}(c_{\mathcal{P}})$.
			Donc par hypothèse de récurrence on a $\ldoubleturn_{\mathcal{S}} \mathcal{P}(c_{\mathcal{P}})$,  et par définition de la satisfaction pour $\exists$, on en déduit $\ldoubleturn_{\mathcal{S}} \exists x\ \mathcal{P}(x)$, c'est-à-dire $\ldoubleturn_{\mathcal{S}} \mathcal{Q}$.
			
			
			\item Si $\ldoubleturn_{\mathcal{S}} \mathcal{Q}$, alors $\ldoubleturn_{\mathcal{S}} \exists x\ \mathcal{P}(x)$, donc 
			il existe un terme fermé $t$, dont l'interprétation est $t^{\mathcal{S}} \in E$, tel que $\ldoubleturn_{\mathcal{S}} \mathcal{P}(t)$.			
			Par hypothèse de récurrence on a $\mathcal{T} \lturn \mathcal{P}(t)$. Ainsi par la règle d'introduction de $\exists$ on a bien $\mathcal{T} \lturn \exists x\ \mathcal{P}(x)$ donc $\mathcal{T} \lturn \mathcal{Q}$.			
		\end{itemize}
		
	\end{itemize}
\end{proof}


%%%%%%%%%%%%%%%%%%%%%%%%%%%%%%%%%%%%%%%%%%%%%%%%%%%%%%%%%%%%%%%%%%%%%
\section{Construction d'une théorie complète}
\label{sec:complete}

Expliquons comment rendre une théorie complète en ajoutant des formules.

%--------------------------------------------------------------------
\subsection{Construction}

Soit $\mathcal{T}$ une théorie cohérente sur un langage $\mathcal{L}$. Nous allons construire une théorie $\mathcal{T}'$, sur le même langage $\mathcal{L}$, telle que $\mathcal{T} \subset \mathcal{T}'$ et $\mathcal{T}'$ est complète.
On va supposer que le langage $\mathcal{L}$ est dénombrable, donc l'ensemble des formules l'est aussi. Considérons une énumération $\mathcal{P}_0, \mathcal{P}_1, \ldots$ de toutes les formules fermées de $\mathcal{L}$.
On construit $\mathcal{T}'$ comme la limite de $(\mathcal{T}_n)_{n\ge0}$ définie par récurrence :
\begin{itemize}
	\item $\mathcal{T}_0 = \mathcal{T}$
	\item 
	\begin{itemize}
		\item $\mathcal{T}_{n+1} = \mathcal{T}_n \cup \{\mathcal{P}_n\}$ si cette union est cohérente,
		\item $\mathcal{T}_{n+1} = \mathcal{T}_n \cup \{\lnot \mathcal{P}_n\}$ sinon.
	\end{itemize}
\end{itemize}
Enfin on pose $\mathcal{T}' = \bigcup_{n \ge 0} \mathcal{T}_n$.

\begin{proposition}
	Chaque $\mathcal{T}_n$ est cohérente et $\mathcal{T}'$ est complète.
\end{proposition}

%--------------------------------------------------------------------
\subsection{Démonstration}

\begin{proof}
	~
	\begin{itemize}
		
		\item Montrons par récurrence que les $\mathcal{T}_n$ sont cohérentes.
		C'est vrai pour $\mathcal{T}_0 = \mathcal{T}$ par hypothèse.
		Supposons que ce soit vrai pour $\mathcal{T}_n$.
		Dans le premier cas, $\mathcal{T}_{n+1} = \mathcal{T}_n \cup \{\mathcal{P}_n\}$ est cohérente par construction.
		Dans le second cas, $\mathcal{T}_n \cup \{\mathcal{P}_n\}$ n'est pas cohérente, donc $\mathcal{T}_n \lturn \lnot \mathcal{P}_n$ (Proposition \ref{prop:noncoherente}). Si $\mathcal{T}_{n+1} = \mathcal{T}_n \cup \{\lnot \mathcal{P}_n\}$ n'était pas cohérente non plus, alors $\mathcal{T}_n \lturn \lnot(\lnot \mathcal{P}_n)$ donc $\mathcal{T}_n \lturn \mathcal{P}_n$.
		On aurait alors $\mathcal{T}_n \lturn \lnot \mathcal{P}_n$ et $\mathcal{T}_n \lturn \mathcal{P}_n$, ce qui contredit l'hypothèse de récurrence que $\mathcal{T}_n$ est cohérente.
		Ainsi, dans tous les cas, $\mathcal{T}_{n+1}$ est cohérente.
		
		
		\item $\mathcal{T}'$ est cohérente. En effet, sinon on aurait $\mathcal{T}' \lturn \lantitauto$. Comme une preuve est constituée d'un nombre fini d'étapes, donc d'un nombre fini de formules, ces formules sont incluses dans un certain $\mathcal{T}_n$, pour $n$ assez grand. On aurait alors $\mathcal{T}_n \lturn \lantitauto$, ce qui contredit la cohérence de $\mathcal{T}_n$.
		
		\item $\mathcal{T}'$ est complète. En effet, soit $\mathcal{Q}$ une formule fermée de $\mathcal{L}$. C'est l'une des $\mathcal{P}_n$. Donc, $\mathcal{P}_n \in \mathcal{T}_{n+1}$ ou $\lnot \mathcal{P}_n \in \mathcal{T}_{n+1}$. Ainsi $\mathcal{Q}$ ou $\lnot \mathcal{Q}$ est un élément de $\mathcal{T}'$. Par conséquent, $\mathcal{T}' \lturn \mathcal{Q}$ ou $\mathcal{T}' \lturn \lnot \mathcal{Q}$.
	\end{itemize}
\end{proof}

Remarque utile pour la suite : si $\mathcal{T}$ admet des témoins de Henkin, alors $\mathcal{T}'$ aussi. En effet \og{}$(\exists x\ \mathcal{P}(x)) \limplies \mathcal{P}(c_{\mathcal{P}})$\fg{} est dans $\mathcal{T}$ donc dans $\mathcal{T}'$.


%%%%%%%%%%%%%%%%%%%%%%%%%%%%%%%%%%%%%%%%%%%%%%%%%%%%%%%%%%%%%%%%%%%%%
\section{Construction des témoins}
\label{sec:temoins}

Partons d'une théorie $\mathcal{T}$ d'un langage $\mathcal{L}$. Nous allons étendre le langage afin d'obtenir une nouvelle théorie $\mathcal{T}'$ sur un langage $\mathcal{L}'$ qui admet des témoins de Henkin, c'est-à-dire :
\mycenterline{\og{}$\exists x\ \mathcal{P}(x) \ \ \limplies \ \ \mathcal{P}(c_{\mathcal{P}})$\fg{} \quad est une formule de $\mathcal{T}'$.}
 En plus $\mathcal{T}'$ doit être cohérente, comme $\mathcal{T}$.

%--------------------------------------------------------------------
\subsection{Extension du langage et de la théorie}

$\mathcal{L}'$ sera construit à partir de $\mathcal{L}$ en ajoutant de nouveaux symboles de constantes $c_0$, $c_1$, \ldots{}
Avant de donner la preuve entière, voici à quoi correspond l'étape fondamentale.


Soit $\mathcal{T}$ une théorie cohérente d'un langage $\mathcal{L}$. Soit $\mathcal{P}$ une formule de $\mathcal{L}$ ayant $x$ pour seule variable libre. Soit $c_{\mathcal{P}}$ un nouveau symbole de constante (qui n'est pas dans $\mathcal{L}$). On note $\mathcal{L}' = \mathcal{L} \cup \{c_{\mathcal{P}}\}$.

\medskip

\textbf{Affirmation.} $\mathcal{T}' = \mathcal{T} \cup \{\exists x\ \mathcal{P}(x) \limplies \mathcal{P}(c_{\mathcal{P}})\}$ est encore cohérente.

On a par construction un témoin de Henkin (juste pour la formule $\mathcal{P}$) et on a conservé la cohérence.

\begin{proof}
	Par l'absurde, si $\mathcal{T}'$ n'est pas cohérente :
	\begin{align*}
		& \mathcal{T} \cup \{\exists x\ \mathcal{P}(x) \limplies \mathcal{P}(c_{\mathcal{P}})\} \lturn \lantitauto \\
		\text{donc } \quad & \mathcal{T} \lturn \lnot (\exists x\ \mathcal{P}(x) \limplies \mathcal{P}(c_{\mathcal{P}})) \\
		\text{donc } \quad & \mathcal{T} \lturn \exists x\ \mathcal{P}(x) \ \land \ \lnot \mathcal{P}(c_{\mathcal{P}}) \qquad \text{car } \lnot(\mathcal{A} \limplies \mathcal{B}) \equiv \mathcal{A} \land \lnot \mathcal{B} \\
		\text{donc } \quad & \mathcal{T} \lturn \exists x\ \mathcal{P}(x)
		\ \text{ et  } \ \mathcal{T} \lturn \lnot \mathcal{P}(c_{\mathcal{P}}) \\		
		\text{donc } \quad & \mathcal{T} \lturn \exists x\ \mathcal{P}(x)
\ \text{ et  } \ \mathcal{T} \lturn \forall x \ \lnot \mathcal{P}(x) \qquad \text{par la règle d'introduction de $\forall$}\\
		\text{donc } \quad & \mathcal{T} \lturn \exists x\ \mathcal{P}(x)
\ \text{ et  } \ \mathcal{T} \lturn \lnot\big(\exists x \ \mathcal{P}(x)\big) \qquad \text{par la négation de $\forall$}\\
		\text{donc } \quad & \mathcal{T} \lturn \mathcal{Q}
\ \text{ et  } \ \mathcal{T} \lturn \lnot \mathcal{Q} \qquad \text{pour $\mathcal{Q} = \exists x \ \mathcal{P}(x)$}
	\end{align*}
	Ce qui contredit $\mathcal{T}$ cohérente. Conclusion : $\mathcal{T}'$ est cohérente.
	Noter l'introduction de $\forall$, afin de passer de \og{}$\lnot \mathcal{P}(c_{\mathcal{P}})$\fg{} à \og{}$\forall x \lnot \mathcal{P}(x)$\fg{} qui est légitime car la constante $c_{\mathcal{P}}$ n'apparaît pas dans $\mathcal{T}$. En effet, $\mathcal{T}$ sont des formules de $\mathcal{L}$, qui ne contient pas la constante $c_{\mathcal{P}}$.
\end{proof}

On pourrait penser faire ceci à chaque formule $\mathcal{P}_0$, $\mathcal{P}_1$,\ldots{} de $\mathcal{L}$, mais ce sera un peu plus délicat que cela. En effet à chaque étape on augmente le langage, donc il y a de plus en plus de formules ce qui complique la récurrence.

%--------------------------------------------------------------------
\subsection{Construction}

Soit $\mathcal{T}$ une théorie cohérente d'un langage $\mathcal{L}$.
On va construire par récurrence une suite de théories $(\mathcal{T}_n)_{n \ge 0}$ et une suite de langages $(\mathcal{L}_n)_{n \ge 0}$ de la façon suivante :
\begin{itemize}
	\item $\mathcal{L}_0 = \mathcal{L}$ et $\mathcal{T}_0 = \mathcal{T}$,
	\item puis pour $n \ge 0$ :
	\begin{itemize}
		\item $\mathcal{L}_{n+1} = \mathcal{L}_n \cup \big\{c_{\mathcal{P}}\big\}_{\mathcal{P}}$ où l'union porte sur toutes les formules $\mathcal{P}$ de $\mathcal{L}_n$ ayant une seule variable libre,
		\item $\mathcal{T}_{n+1} = \mathcal{T}_n \cup \big\{\exists x\ \mathcal{P}(x) \limplies \mathcal{P}(c_{\mathcal{P}})\big\}_{\mathcal{P}}$.
	\end{itemize}
\end{itemize}
Enfin, on pose $\mathcal{L}' = \bigcup_{n \ge 0} \mathcal{L}_n$ et $\mathcal{T}' = \bigcup_{n \ge 0} \mathcal{T}_n$.

\medskip

\textbf{Affirmation.} $\mathcal{T}'$ admet des témoins de Henkin.

\begin{proof}
	Soit $\mathcal{P}$ une formule de $\mathcal{L}'$ ayant une seule variable libre $x$. Elle contient un nombre fini de symboles de $\mathcal{L}'$, donc pour un $n$ suffisamment grand, $\mathcal{P}$ est une formule de $\mathcal{L}_n$. Donc $\exists x\ \mathcal{P}(x) \limplies \mathcal{P}(c_{\mathcal{P}})$ est une formule de $\mathcal{T}_{n+1}$ donc de $\mathcal{T}'$.
\end{proof}

\medskip

\textbf{Affirmation.} $\mathcal{T}'$ est cohérente.

\begin{proof}
	On va montrer que chaque $\mathcal{T}_n$ est cohérente. Cela prouvera que $\mathcal{T}'$ est cohérente. En effet si on avait $\mathcal{T}' \lturn \lantitauto$, alors une preuve de ceci ne fait intervenir qu'un nombre fini de formules de $\mathcal{T}'$, qui sont donc toutes incluses dans un $\mathcal{T}_n$ pour un $n$ assez grand. Pour ce $n$, on a $\mathcal{T}_n \lturn \lantitauto$ ce qui contredit $\mathcal{T}_n$ cohérente.
	
	La preuve que les $\mathcal{T}_n$ sont cohérentes se fait par récurrence, sur le même principe que l'étape fondamentale expliquée pour une seule formule au début de cette section.
	
	\medskip
	
	Montrons par récurrence que $\mathcal{T}_n$ (pour $n \ge 0$) est cohérente.
	Pour $n=0$, $\mathcal{T}_0 = \mathcal{T}$, alors $\mathcal{T}_0$ est cohérente par hypothèse.
	
	Supposons $\mathcal{T}_n$ cohérente et par l'absurde que $\mathcal{T}_{n+1}$ ne l'est pas :
	\begin{align*}
		& \mathcal{T}_n \cup \bigcup_{i=1}^k \{\exists x\ \mathcal{P}_i(x) \limplies \mathcal{P}_i(c_{\mathcal{P}_i})\} \lturn \lantitauto \qquad \text{(avec un nombre fini de $\mathcal{P}_i$ impliquées dans la preuve)} \\
		\text{donc } \quad & \mathcal{T}_n \cup  \big\{\bigwedge_{i=1}^k \big(\exists x\ \mathcal{P}_i(x) \limplies \mathcal{P}_i(c_{\mathcal{P}_i})\big)\big\} \lturn \lantitauto \\
		\text{donc } \quad & \mathcal{T}_n \lturn \lnot \bigwedge_{i=1}^k \big(\exists x\ \mathcal{P}_i(x) \limplies \mathcal{P}_i(c_{\mathcal{P}_i})\big) \\
		\text{donc } \quad & \mathcal{T}_n \lturn \bigvee_{i=1}^k \lnot \big(\exists x\ \mathcal{P}_i(x) \limplies \mathcal{P}_i(c_{\mathcal{P}_i})\big) \\
		\text{donc } \quad & \mathcal{T}_n \lturn \lnot \big(\exists x\ \mathcal{P}_j(x) \limplies \mathcal{P}_j(c_{\mathcal{P}_j})\big) \qquad \text{pour un certain } j \text{ (par le Lemme \ref{lem:compbis})} \\
		\text{donc } \quad & \mathcal{T}_n \lturn \exists x\ \mathcal{P}_j(x) \ \text{ et } \ \mathcal{T}_n \lturn \lnot \mathcal{P}_j(c_{\mathcal{P}_j}) \\
		\text{donc } \quad & \mathcal{T}_n \lturn \exists x\ \mathcal{P}_j(x) \ \text{ et } \ \mathcal{T}_n \lturn \forall x\ \lnot \mathcal{P}_j(x) \qquad \text{par la règle d'introduction de $\forall$} \\
		\text{donc } \quad & \mathcal{T}_n \lturn \exists x\ \mathcal{P}_j(x) \ \text{ et } \ \mathcal{T}_n \lturn \lnot \big(\exists x\ \mathcal{P}_j(x)\big)
	\end{align*}
	Donc $\mathcal{T}_n$ est non cohérente. C'est contradictoire avec l'hypothèse de récurrence.
	Conclusion : $\mathcal{T}_{n+1}$ est cohérente.
	
	Noter encore une fois l'introduction de $\forall$, afin de passer de \og{}$\lnot \mathcal{P}_j(c_{\mathcal{P}_j})$\fg{} à \og{}$\forall x \ \lnot \mathcal{P}_j(x)$\fg{} qui est légitime car la constante $c_{\mathcal{P}_j}$ n'apparaît pas dans $\mathcal{T}_n$.
\end{proof}

%--------------------------------------------------------------------
\subsection{Preuve complète du théorème}

Nous avons donné les détails de la preuve du théorème de complétude mais un peu dans le désordre. Voici le plan de la preuve remis dans le bon sens.

\begin{itemize}
	\item \textbf{Hypothèses.} $\mathcal{T}$ théorie cohérente, $\mathcal{L}$ langage dénombrable.
	
	\item \textbf{Étape 1 : Ajouter les témoins de Henkin.} 
	On renvoie au début de cette section \ref{sec:temoins}.
	On obtient une théorie $\mathcal{T}'$ sur un nouveau langage $\mathcal{L}'$ avec $\mathcal{T} \subset \mathcal{T}'$ et $\mathcal{L} \subset \mathcal{L}'$.
	$\mathcal{T}'$ admet des témoins de Henkin et reste cohérente.
	$\mathcal{L}'$ est obtenu à partir de $\mathcal{L}$ en ajoutant un nombre dénombrable de constantes $c_{\mathcal{P}}$, ainsi $\mathcal{L}'$ reste dénombrable.
	
	\item \textbf{Étape 2 : Rendre la théorie complète.} 
	Voir la section \ref{sec:complete}.
	À partir de $\mathcal{T}'$ et $\mathcal{L}'$, on obtient une théorie $\mathcal{T}''$ complète avec $\mathcal{T}' \subset \mathcal{T}''$.	
	De plus on préserve les témoins de Henkin (remarque à la fin de la section \ref{sec:complete}). Le langage est inchangé, on note $\mathcal{L}'' = \mathcal{L}'$.
	
	\item \textbf{Étape 3 : Construction du modèle.}
	C'est la section \ref{sec:modele} appliquée à $\mathcal{T}''$ et $\mathcal{L}''$.
	On obtient un modèle $\mathcal{S}''$ de $\mathcal{T}''$ sur $\mathcal{L}''$.
	Ce modèle se restreint naturellement en un modèle $\mathcal{S}$ de $\mathcal{T}$ sur $\mathcal{L}$.
	En effet, $\mathcal{L}''$ est obtenu à partir de $\mathcal{L}$ en ajoutant des constantes $c_{\mathcal{P}}$.
	On considère la structure $\mathcal{S}$ identique à $\mathcal{S}''$ sauf que l'on ne donne pas d'interprétation de ces constantes (qui n'existent pas dans $\mathcal{L}$).
	Soit $\mathcal{P}$ une formule de $\mathcal{T}$ sur $\mathcal{L}$. 
    C'est donc aussi une formule de $\mathcal{T}''$ sur $\mathcal{L}''$.
    Comme $\mathcal{S}''$ est un modèle de $\mathcal{T}''$, alors on a $\ldoubleturn_{\mathcal{S}''} \mathcal{P}$.
	Mais comme $\mathcal{P}$ est une formule de $\mathcal{L}$, elle ne contient pas de symboles de ces constantes supplémentaires, donc on a aussi $\ldoubleturn_{\mathcal{S}} \mathcal{P}$.
	Ainsi, $\mathcal{S}$ est un modèle pour $\mathcal{T}$.

\end{itemize}


\bigskip 
\bigskip


Ainsi se termine la preuve du théorème de complétude de Gödel.
Ce théorème était le sujet de sa thèse, terminée en 1929 à l'âge de 23 ans et publiée l'année d'après.
Ainsi se termine aussi cette partie sur la logique du premier ordre.
On continue dans la partie suivante avec ce qui a rendu Gödel célèbre dans le monde entier bien au-delà des mathématiques : le théorème d'incomplétude.

%%%%%%%%%%%%%%%%%%%%%%%%%%%%%%%%%%%%%%%%%%%%%%%%%%%%%%%%%%%%%%%%%%%%%
\section*{Exercices}

\subsection*{Le théorème de complétude}

\begin{exercicecours}
	On rappelle les théorèmes A et B :
	\begin{itemize}
		\item Théorème A. Si $\mathcal{T} \ldoubleturn \mathcal{Q}$, alors $\mathcal{T} \lturn \mathcal{Q}$.
		
		\item Théorème B. Si $\mathcal{T}$ est cohérente, alors $\mathcal{T}$ admet un modèle.		
	\end{itemize}
	
	On admet le théorème A. En déduire le théorème B.	
\end{exercicecours}


\begin{exercicecours}
	Une théorie $\mathcal{T}$ (qui peut contenir une infinité de formules) est dite \emph{finiment cohérente} si pour tout sous-ensemble fini $\Gamma \subset \mathcal{T}$, $\Gamma$ est cohérente.
	
	Montrer qu'une théorie $\mathcal{T}$ est cohérente si et seulement si elle est finiment cohérente.
\end{exercicecours}



\subsection*{Théorie complète}

\begin{exercicecours}
	Soit $\mathcal{L}$ un langage du premier ordre et soit $\mathcal{S}$ une structure pour ce langage. 
	On définit la \emph{théorie} de $\mathcal{S}$, notée $\mathcal{T}(\mathcal{S})$, comme l'ensemble de toutes les formules fermées $\mathcal{P}$ sur $\mathcal{L}$ qui sont vraies dans $\mathcal{S}$ :
	\[ \mathcal{T}(\mathcal{S}) = \big\{ \mathcal{P} \text{ formule fermée} \mid \ \  \ldoubleturn_{\mathcal{S}} \mathcal{P} \big\} \]

	Montrer que $\mathcal{T}(\mathcal{S})$ est une théorie complète.
\end{exercicecours}


\begin{exercicecours}
	Soit $\mathcal{T}$ une théorie cohérente. Montrer l'équivalence entre :
	\begin{itemize}
		\item[(i)] $\mathcal{T}$ est complète,
		
		\item[(ii)] pour toute formule fermée $\mathcal{Q}$, si $\mathcal{T} \cup \{\mathcal{Q}\}$ est cohérente, alors $\mathcal{T} \lturn \mathcal{Q}$.
	\end{itemize}
\end{exercicecours}


\subsection*{Construction d'un modèle à partir d'une théorie complète}

\begin{exercicecours}
	On considère un langage $\mathcal{L}$ sans symbole d'égalité, contenant une constante $0$, un symbole de fonction unaire $s$, et un prédicat unaire $P$.
	Soit $\mathcal{T}$ une théorie complète sur $\mathcal{L}$ qui contient notamment les deux axiomes (fomules) suivants :
	\begin{itemize}
		\item $\mathcal{A}_1 : P(0)$
		\item $\mathcal{A}_2 : \forall x \ \  \big( P(x) \  \lequiv \ \lnot P(s(x)) \big)$
	\end{itemize}
	
	\begin{enumerate}
		\item Construire le modèle $\mathcal{S}$ canonique associé à $\mathcal{T}$ (comme dans la Section \ref{ssec:modele}).
		
		\item Construire un modèle $\mathcal{S}'$ avec la structure de domaine $\Nn$, 
		où $s(x)$ correspond à $x+1$ (le successeur) et expliquer à quoi correspond $P(x)$.	
	\end{enumerate}

\end{exercicecours}


\begin{exercicecours}
Compléter la preuve du Lemme \ref{lem:satisfaction} pour les connecteurs manquants, ainsi que pour le quantificateur $\exists$, c'est-à-dire pour les formules :
$\mathcal{Q} = \mathcal{P}_1 \lor \mathcal{P}_2$, 
$\mathcal{Q} = \mathcal{P}_1 \limplies \mathcal{P}_2$, 
$\mathcal{Q} = \mathcal{P}_1 \lequiv \mathcal{P}_2$, 
$\mathcal{Q} = \exists x \ \mathcal{P}(x)$. 
\end{exercicecours}


\subsection*{Construction d'une théorie complète}


\begin{exercicecours}
	On se place dans le cadre plus simple de la logique propositionnelle.
	Soit un langage $\mathcal{L}$ contenant seulement les deux variables propositionnelles $A$ et $B$.
	
	\begin{enumerate}
	\item Justifier que toute formule est une formule fermée. Donner des exemples de formules.
	
	\item Justifier que construire un modèle d'une théorie $\mathcal{T}$, c'est attribuer une valeur V/F à $A$ et une valeur V/F à $B$ qui rendent toutes les formules de $\mathcal{T}$ vraies.
	
	\item Justifier que si $\mathcal{T}$ admet un modèle, alors $\mathcal{T}$ est cohérente.
	\end{enumerate}
	
	
	On veut maintenant construire (le début d') une théorie complète à partir
	de $\mathcal{T}_0 = \{ A \}$.
	On considère le choix suivant pour l'énumération $(\mathcal{P}_n)_{n\ge0}$ des formules de $\mathcal{L}$ :
	\[
	B, \quad A \lor B, \quad \lnot A, \quad A \limplies B, \quad \ldots
	\]
	Pour $n \ge 0$ on définit :
	\begin{itemize}
		\item $\mathcal{T}_{n+1} = \mathcal{T}_n \cup \{\mathcal{P}_n\}$ si cette union est cohérente,
		\item $\mathcal{T}_{n+1} = \mathcal{T}_n \cup \{\lnot \mathcal{P}_n\}$ sinon.
	\end{itemize}
	
	\begin{enumerate}
	\setcounter{enumi}{3}
	\item Justifier que $\mathcal{T}_0$ est cohérente.
	
	\item Construire $\mathcal{T}_1$, \ldots, $\mathcal{T}_4$.
	
	\item La théorie complète $\mathcal{T} = \bigcup_{n\ge0} \mathcal{T}_n$  prouvera-t-elle $A \land B$ ?
	
	\item On repart de $\mathcal{T}'_0 = \{ A \}$, mais cette fois on considère l'énumération :
	\[
	\lnot B, \quad A \lor B, \quad \lnot A, \quad A \limplies B, \quad \ldots
	\]	
	Construire $\mathcal{T}'_1$, \ldots, $\mathcal{T}'_4$.
	
	\item La théorie complète $\mathcal{T}' = \bigcup_{n\ge0} \mathcal{T}'_n$
	sera-t-elle identique à $\mathcal{T}$ ?
	\end{enumerate}
\end{exercicecours}


\subsection*{Construction des témoins}

\begin{exercicecours}
	On considère un langage initial $\mathcal{L}_0$ contenant un unique prédicat unaire $P(x)$, sans constante ni fonction. 
	Soit la théorie $\mathcal{T}_0$ composée des deux axiomes suivants :
	\[ 
	\mathcal{A}_1 : \exists x\ P(x) \qquad \text{et} \qquad \mathcal{A}_2 : \exists x\ \lnot P(x) 
	\]
	
	\begin{enumerate}
		\item Justifier que $\mathcal{T}_0$ est cohérente (par exemple en construire un modèle).
		
		\item Le langage $\mathcal{L}_0$ possède-t-il des termes clos ? La théorie $\mathcal{T}_0$ admet-elle des témoins de Henkin sur $\mathcal{L}_0$ ?
		
		\item On applique la méthode de Henkin pour les deux formules $P(x)$ et $\lnot P(x)$. On note alors $\mathcal{L}_1 = \mathcal{L}_0 \cup \{c_1, c_2\}$ le langage étendu par deux constantes $c_1$ et $c_2$.
		Quels axiomes de Henkin a-t-on ajouté ? Que vaut alors la théorie $\mathcal{T}_1$ ?
		
		\item Montrer que $\mathcal{T}_1$ prouve $P(c_1) \land (\lnot P(c_2))$.
		En déduire que dans tout modèle de $\mathcal{T}_1$, l'interprétation de $c_1$ et l'interprétation de $c_2$ sont nécessairement deux éléments distincts du domaine. 
	\end{enumerate}
\end{exercicecours}


\begin{exercicecours}
	Soit $\mathcal{L}'$ une extension d'un langage $\mathcal{L}$.
	Soit $\mathcal{T}$ une théorie sur le langage $\mathcal{L}'$, qui est complète et qui admet des témoins de Henkin.
	Soit $\mathcal{P}(x)$ une formule du langage initial $\mathcal{L}$ (ayant $x$ pour seule variable libre).
	
	Montrer que si $\mathcal{T}$ ne prouve pas $\forall x\ \mathcal{P}(x)$, alors il existe une constante de Henkin $c$ du langage $\mathcal{L}'$ telle que :
	\[ \mathcal{T} \quad \lturn \quad \lnot \mathcal{P}(c). \]
	
	Comment appelle-t-on $c$ en langage courant ?
\end{exercicecours}


\end{document}